\section{Marco metodológico}

\begin{frame}{Población y muestreo}
	\centering
	\begin{itemize}
		\item Autómatas celulares finitos definidos en un espacio bidimensional con una cuadrícula regular.
		\item Condiciones de frontera.
		\item Conjunto finito de estados.
		\item Cualquier función de transición que mapee de $\mathcal{R}^n \to \mathcal{R}$
	\end{itemize}
	Se selecciona un subconjunto mediante un muestreo no probabilístico de tipo intencional.
\end{frame}

\begin{frame}{Investigación}
	\begin{itemize}
		\item Tipo aplicada.
		\item Diseño cuasiexperimental.
		\item Enfoque cuantitativo.
	\end{itemize}
\end{frame}

\begin{frame}{Procedimiento}
	\begin{itemize}
		\item Diseño, desarrollo e implementación del simulador de autómatas celulares.
		
		\item Diseño, desarrollo e integración del módulo de red neuronal artificial, encargado de aproximar la función de transición.
		
		\item Imitación de patrones conocidos, a fin de validar el desempeño del modelo neuronal frente a sistemas con comportamiento bien caracterizado.
		
		\item Simulación de otros sistemas complejos, evaluando la capacidad generalizadora de la aproximación propuesta.
		
		
	\end{itemize}
\end{frame}
